\section*{Executive Summary}
\addcontentsline{toc}{section}{Executive Summary}

\textbf{To:} The Office of the Mayor and the Sangguniang Bayan of Maigo \\
\textbf{Subject:} Institutionalizing Dual-Use ICT Facilities via Public-Private Partnership \\
\textbf{Framework:} RA 11966 (PPP Code of the Philippines)

\vspace{1em}
\noindent \textbf{The Core Problem: Infrastructure Scarcity vs. Asset Underutilization} \\
The Municipality of Maigo faces a critical bottleneck in its digital upskilling initiatives. Key educational institutions, including MSU-MCEST and local TESDA branches, lack the capital expenditure (CapEx) required to build and maintain high-end computer laboratories for Technical-Vocational Education and Training (TVET). Consequently, residents are deprived of advanced ICT certifications. Simultaneously, local Micro, Small, and Medium Enterprises (MSMEs) operating internet cafés possess high-performance hardware that sits idle and depreciates during daytime academic hours due to shifts in mobile internet usage.

\vspace{1em}
\noindent \textbf{The Proposed Solution: The "Dual-Use" PPP Model} \\
This policy research proposes the institutionalization of a service-contracting model under the 2023 PPP Code. Instead of purchasing redundant hardware, the Local Government Unit (LGU) will lease accredited private cybercafés during standard academic hours (8:00 AM -- 5:00 PM) to serve as public training micro-laboratories. 

To ensure a secure and distraction-free learning environment, the policy mandates a \textbf{Diskless (PXE-Boot) Multiple Image architecture}. This technical requirement allows the hardware to boot into a sterile, government-controlled ``Educational Image'' (containing only TVET software and blocking all games/social media) during the day, before reverting to standard commercial operations at night.

\vspace{1em}
\noindent \textbf{Strategic Benefits to the Municipality of Maigo:}
\begin{itemize}
    \item \textbf{Zero-CapEx Upgrades:} The LGU and public institutions bypass the multi-million peso costs of procuring hardware, facility construction, and IT maintenance.
    \item \textbf{MSME Economic Resilience:} Struggling local internet cafés receive a stable, institutionalized revenue stream through daytime government leasing, preventing business closures.
    \item \textbf{Bridging the Digital Divide:} Out-of-school youth and students gain immediate access to industry-standard hardware for advanced certifications (e.g., AutoCAD, CSS NC II, Programming).
\end{itemize}

\vspace{1em}
\noindent \textbf{Actionable Output} \\
This document outlines the comparative policy alternatives and provides a draft \textbf{Municipal Ordinance} and \textbf{Service Level Agreement (SLA)} to legally and technically operationalize this framework.